\section{Detecting integral membership from the primitive levels of an arc}

% BEGIN D82P1
\paragraph{D82P1. The filtered Taylor test for the physical Wach lattice.}
Keep J82F1--J82F2's notation, on either tested lattice of the
actual unramified row B. Put $N=N_c$ and
$W=\Lambda^{-1}A\varphi^\ell N_{c-\ell}$.
For each $0\le s<\ell$ the level filtration in $D_c$ is
\[
        F_s^u=\varphi_{\rm cris}^s
          (\mathrm{Fil}^uD_{{\rm dR},c-s}),
\]
using the unramified crystalline/de Rham identification at
$c-s$. If $x\in W$, then $x\in N$ if and only if at every
primitive root of order $p^{s+1}$, for every $s<\ell$, the
coefficient of $Y^k$ in $x$ belongs to $F_s^{-k}$ for each
$k\in\mathbb Z$. One may take a common finite splitting field
for the coefficient tests.

At a primitive $p$th root the rational Taylor lattice of $N_c$
is exactly this filtered Laurent lattice by W79H2--W79H3. An
invariant Wach basis also gives a basis over the original series
ring by L81S2, so the span in the Taylor ring can be computed in
that basis. At a primitive $p^{s+1}$th root with $s>0$ the
height factor $Q$ is a Taylor unit. Both original height
inclusions identify the Taylor lattice there with the Frobenius
translate of the preceding letter at the $p$th power of the
root. The substitution $Y\mapsto pY$ is an automorphism over
the coefficient field. After $s$ steps this proves the displayed
filtration formula for the span at that root. All choices of
the final primitive $p$th root give the same filtration because
there is a single embedding of the unramified base extending
the letter. This also gives the source's rational local
submodule $\mathrm{Nl}_s$, whose degree $k$ condition is
membership in $F_s^{-k}$ on each branch.

To recover integral membership, a product
$D=\prod_{s<\ell}P_s^{m_s}$ puts $x$ in $N$:
clear $\Lambda$ and successively use the lower Wach height
inclusions to put the generating columns of the Frobenius span
in $N$. In the free invariant basis of $N$ write the coordinates
of $x$ as $a_m/D$ with $a_m\in A$. Its Taylor expansion lies
in the span just computed precisely when all these coordinates
are regular. Indeed the expanded basis is independent over
the Laurent field and spans this Taylor lattice. Regularity at
every branch forces the Taylor order of $a_m$ at each zero
of $P_s$ to be at least $m_s$. The map from a Taylor parameter
at each nonzero tame root in $v$ to $Y$ preserves order:
$P_s$ has a simple zero there and its expansion has a simple
zero in $Y$, as in J82F2. Divide $a_m$ by the monic
distinguished polynomial associated to $D$ over $O$. Its
polynomial remainder has all the roots just described with
their multiplicities, so vanishes by its degree bound, even
if the roots were split over an extension. Thus $D\mid a_m$
already in $A$. This proves the reverse implication integrally;
the forward implication follows from the Taylor lattice calculation.
It proves the native $\mathrm{Mem}$ on the full $W$.
% END D82P1

% BEGIN D82P2
\paragraph{D82P2. Clearing the common denominator and detecting multiplication by the tame parameter.}
Put $\Theta=\prod_{s<\ell}P_s^g$. For a thin, up and down
position let $h_s$ denote respectively $G$, $b_s+S$, and
$\eta$, using the width $S$ of the marked jump band. These are
the upper jump bounds with $H_s=g+h_s$. The upper height
inclusion of W79H4 at letter $c-s$ is
$P_0^{h_s}N_{c-s}\subset A\varphi N_{c-s-1}$ in invariant
coordinates. Successively applying it, and applying
$\varphi^s$ to the $s$th scalar, gives
\[
          (\Lambda/\Theta)N
              \ \subseteq\ A\varphi^\ell N_{c-\ell}.
\]
For every $x$ in the ambient \'etale module with $\Theta x\in N$
this implies $\Lambda x\in A\varphi^\ell N_{c-\ell}$.
In particular $\Theta^{-1}N\subset W$ and $N\subset W$.

For $x\in W$, D82P1 applied to $\Theta x\in W$ gives
\[
 \Theta x\in N\quad\Longleftrightarrow\quad
 P_s^g\,\mathrm{loc}_s(x)\in\mathrm{Nl}_s
                              \quad\hbox{for every }s<\ell .
\]
The remaining factors of $\Theta$ are Taylor units at all roots
at level $s$, and any Taylor unit and its inverse preserve the
filtered Laurent lattice. This proves both clauses of
$\mathrm{DivisibleIff}(g)$.
If $x\in vN$ then $x\in vW$ and its localizations belong
to $\mathrm{Nl}_s$. Conversely, write $x=vy$ with $y\in W$
and assume these localizations of $x$ belong to $\mathrm{Nl}_s$.
At all primitive roots $v$ is a unit. Therefore the same tests
hold for $y$ and D82P1 gives $y\in N$. This proves
$\mathrm{XMulIff}$ with the native variable $v$.
% END D82P2

% BEGIN D82P3
\paragraph{D82P3. Rounded filtration bounds and containment of the frame error.}
At a fixed level choose a rational basis adapted to $F_s$.
A basis vector with jump $a$ contributes the Laurent line
$Y^{-a}L[[Y]]$ to $\mathrm{Nl}_s$ (with scalars extended
when necessary). On the thin row round every jump to $0$.
On an up row round the jumps in $[b_s,b_s+S]$ to $b_s$ and
the central jumps to $0$. On a down row round the jumps in
$[-b_s-S,-b_s]$ to $-b_s$ and the central jumps to $0$.
There is no ambiguity because $b_s>g\ge\eta$ at a mark.
These give the graded lines of the native $\mathrm{Nb}_s$:
the up subspace is $F_s^{b_s}$ and the down subspace surviving
in degrees $0,\ldots,b_s-1$ is $F_s^{-\eta}$. In every
case $|a-\bar a|\le g$ by I82B3. Thus multiplying either
Laurent lattice by $Y^g$ takes it into the other one, and $P_s$
is $Y$ times a Taylor unit. These are the two inclusions in
$\mathrm{RD}(g)$ on the level module.

By J82F2 the frame error has order at least $d_s-H_s$.
On a thin row $d_s\ge3g$ and $H_s=g+G\le2g$, so this order
is at least $g$. On an up row $d_s=qM>h+2g$ and
$H_s=g+b_s+S\le g+h$, giving order at least $g$.
All true jumps on either of these rows are at least $-g$,
and all rounded jumps there are nonnegative. On a down row
$d_s=qM>h+2g$, $H_s=g+\eta\le2g$, and $h\ge b_s+S$.
The true and rounded jumps there are at least $-(b_s+S)$.
In every case the error order is therefore at least $-a$ and
$-\bar a$ on every line. Hence
$\mathrm{Fl}_s\subseteq\mathrm{Nl}_s\cap\mathrm{Nb}_s$ for
all $s<\ell$. This is the native $\mathrm{P3}$.
% END D82P3
